\documentclass[11pt,a4paper]{article}

% --- Packages ---
\usepackage[utf8]{utf8}
\usepackage[margin=1in]{geometry}
\usepackage{amsmath,amssymb,amsfonts,amsthm}
\usepackage{graphicx}
\usepackage{xcolor}
\usepackage{hyperref}
\usepackage{booktabs}
\usepackage{enumitem}

% --- Hyperlink Setup ---
\hypersetup{
    colorlinks=true,
    linkcolor=blue,
    filecolor=magenta,      
    urlcolor=cyan,
    pdftitle={Academic Paper / Technical Report},
    pdfpagemode=FullName,
}

% --- Theorem Environments ---
\newtheorem{theorem}{Theorem}[section]
\newtheorem{lemma}[theorem]{Lemma}
\newtheorem{definition}{Definition}[section]

\title{\textbf{Title of the Document}}
\author{\textbf{Aarav Shah} \\ \small Department of Computer Science \\ \small \href{mailto:email@example.com}{email@example.com}}
\date{\today}

\begin{document}

\maketitle

\begin{abstract}
This is a clean abstract summarizing the main objectives, methodologies, results, and implications of this paper.
\end{abstract}

\section{Introduction}
Welcome to your LaTeX document. You can cite equations like Equation~\ref{eq:example} or refer to sections easily.

\begin{equation}
\label{eq:example}
E = mc^2
\end{equation}

\section{Methodology}
Describing the primary technical process or algorithm.

\begin{theorem}
For any non-negative integer $n$, the summation of the first $n$ integers is given by $\frac{n(n+1)}{2}$.
\end{theorem}

\section{Results \& Evaluation}
\begin{table}[h]
\centering
\caption{Performance Overview}
\label{tab:results}
\begin{tabular}{@{}lrr@{}}
\toprule
\textbf{Metric} & \textbf{Baseline} & \textbf{Proposed System} \\
\midrule
Accuracy        & 84.5\%            & 96.2\%                   \\
Latency (ms)    & 120ms             & 42ms                     \\
\bottomrule
\end{tabular}
\end{table}

\section{Conclusion}
Final remarks and future work directions.

\end{document}
